\section{Preliminaries}\label{sec:preliminaries}

Let $\lambda_k$ denote $k$-dimensional Lebesgue measure.  We use the
zero-dimensional convention $\lambda_0(\{()\})=1$; this makes the
$N=1$ coordinate sweep a literal specialization of the general formulas.

For an interval $I\subseteq\Theta$, define the coefficient-space root event
\[
S_I
=\{\alpha\in[-R,R]^N:
  \exists\theta\in I,\ \phi_\alpha(\theta)=0\}.
\]
This notation names the event whose volume is swept in the proof and whose
probability appears in the main theorem.  It is defined entirely from the
setting objects; in particular, it introduces no regularity or simple-root
condition.

The normalization $A=(2R)^N\kappa$ is the density cap multiplied by the
volume of the coefficient cube.  It is used only to rewrite
$\kappa(2R)^{N-1}$ as $A/(2R)$ in Theorem~\ref{thm:main}.  No additional helper
constant is needed in the public statement.
